We use SHAP to examine whether a modern image classifier
(ConvNeXt-Tiny, pretrained on ImageNet-1k) attends to the visual features a
human would use to distinguish between visually similar object classes --
ear shape, fur pattern, beak shape -- or relies on incidental cues. We build
a test set of 22 images spanning 10 groups of confusable ImageNet classes
(dog breeds, cat breeds, birds, sharks, big cats, foxes, elephants, and
bears) and compute SHAP attributions for the model's top prediction on each
image. The model misclassifies 5 of 22 images (23\%); four of the five
errors occur between the intended confusable pair. SHAP attribution maps
show that in most of these errors the model does attend to the correct
region of the object (head, face, or body) rather than the background, but
the specific texture cue that separates the two classes -- coat spot
pattern, ear size -- is frequently under-weighted or obscured in the source
image. We discuss this as evidence that errors on confusable classes are
driven less by attending to the wrong part of the image and more by the
salience of the distinguishing texture within the correct region.
